\documentclass[10pt,a4paper]{article}
\newcommand{\ReportShortTitle}{Polar structure and D2 inference}
\usepackage[T1]{fontenc}
\usepackage[utf8]{inputenc}
\usepackage{lmodern}
\usepackage{microtype}
\usepackage[a4paper,margin=24mm,headheight=14pt]{geometry}
\usepackage{amsmath,amssymb,mathtools}
\usepackage{booktabs,tabularx,array}
\usepackage{enumitem}
\usepackage{xcolor}
\usepackage{fancyhdr}
\usepackage{titlesec}
\usepackage[hidelinks]{hyperref}

\definecolor{ink}{HTML}{18212B}
\definecolor{accent}{HTML}{135F73}
\definecolor{muted}{HTML}{5E6A73}
\definecolor{pale}{HTML}{EEF5F7}
\color{ink}
\hypersetup{colorlinks=true,linkcolor=accent,urlcolor=accent}
\setlength{\parindent}{0pt}
\setlength{\parskip}{0.55em}
\setlist{nosep,leftmargin=1.55em}
\setcounter{secnumdepth}{2}
\titleformat{\section}{\Large\bfseries\color{accent}}{\thesection}{0.6em}{}
\titleformat{\subsection}{\large\bfseries\color{ink}}{\thesubsection}{0.55em}{}
\pagestyle{fancy}
\fancyhf{}
\fancyhead[L]{\small ODMA--URA research record}
\fancyhead[R]{\small\color{muted}\ReportShortTitle}
\fancyfoot[C]{\thepage}
\renewcommand{\headrulewidth}{0.3pt}
\newcolumntype{Y}{>{\raggedright\arraybackslash}X}
\newcommand{\R}{\mathbb{R}}
\newcommand{\C}{\mathbb{C}}
\newcommand{\Z}{\mathbb{Z}}
\newcommand{\F}{\mathbb{F}}
\newcommand{\one}{\mathbf{1}}
\newcommand{\given}{\,|\,}
\newcommand{\CN}{\mathcal{CN}}
\DeclareMathOperator{\diag}{diag}
\DeclareMathOperator{\Cov}{Cov}
\DeclareMathOperator{\Var}{Var}
\DeclareMathOperator{\supp}{supp}
\DeclareMathOperator*{\argminop}{arg\,min}
\DeclareMathOperator*{\argmaxop}{arg\,max}
\newcommand{\evidence}[1]{\textcolor{accent}{\textbf{Evidence status:}} #1}
\providecommand{\ReportUpdated}{30 August 2026}
\newcommand{\ProjectReport}[4]{%
  \pagenumbering{roman}
  \begin{titlepage}
    \vspace*{18mm}
    {\large\color{accent}\textbf{ODMA--URA research record}}\\[12mm]
    {\huge\bfseries #1}\\[5mm]
    {\Large\color{muted}#2}\\[18mm]
    \begin{tabularx}{\textwidth}{@{}p{0.16\textwidth}Y@{}}
      Stage & #3\\
      Updated & \ReportUpdated\\
      Status & #4
    \end{tabularx}
    \vfill
    \textcolor{muted}{This report is a chronological research record. Measured results, interpretations, and open
    hypotheses are labelled separately. Detailed numerical ledgers remain in \texttt{results/}.}
  \end{titlepage}
  \pagenumbering{arabic}
}

\newcommand{\sigm}{\operatorname{sigmoid}}
\newcommand{\softplus}{\operatorname{softplus}}
\newcommand{\GF}{\operatorname{GF}}
\newcommand{\E}{\mathbb{E}}
\begin{document}
{\LARGE\bfseries Polar structure and D2 inference}\par
{\large From an explicit message matrix to structured multiuser decoding}\par
{\color{muted}8 October 2026 --- derivation, related work, and Job 034}

\textbf{Question.} Can we preserve the useful joint interference treatment of D2 without keeping one state per
possible message? Polar coding gives a compact description of valid codewords, but does not by itself solve activity,
pattern or multiuser inference. This note locates the overlap and the remaining approximation choices.

\textbf{Status.} Only the symbol-energy and shared-mixer encoder experiments are implemented here. Both retain
the existing, full-message D2 recurrence. A new polar-based receiver is deliberately not selected or implemented.
The evidence motivating this work is in the supervisor report; this note develops the algebra rather than repeating
the result tables. Literature claims are cited; the structural identities below are derived for our model.

\section{One signal model, two descriptions of the encoder}
For $K$ simultaneous transmissions in the real Gaussian MAC,
\begin{equation}
 y=\sum_{k=1}^{K}\phi(w_k)+z=\Phi a^\star+z,
 \qquad z\sim\mathcal N(0,\sigma^2 I_n),\quad
 \Phi=[\phi(w)]_{w\in\{0,1\}^B}\in\R^{n\times M},\quad M=2^B.
\end{equation}
Each column has energy one. The vector $a^\star\in\Z_{\geq0}^M$ counts messages and sums to $K$.
It is binary when messages are distinct. Our comparisons supply $K$ and $\sigma^2$ to the receiver;
with physical $E_b/N_0=\eta$, $\sigma^2=1/(2B\eta)$. No unknown fading is present.

Split the message as $w=(p,v)$, with $B_p$ pattern bits and $B_c=B-B_p$ payload bits. Let
$P=2^{B_p}$ and $V=2^{B_c}$, so $M=PV$. The actual CRC and polar encoder are linear over binary arithmetic:
\begin{equation}
 v\in\GF(2)^{B_c}\ \longmapsto\ c(v)=L v\pmod 2\in\GF(2)^{N_c}
 \ \longmapsto\ b(v)=(-1)^{c(v)}\in\{-1,+1\}^{N_c}.
 \label{eq:polar}
\end{equation}
Here $L$ has shape $N_c\times B_c$ and includes CRC generation, insertion into unfrozen positions and the polar
transform. Binary arithmetic describes the \emph{code}, not the physical channel alphabet. The transmitted waveform is
\begin{equation}
 \boxed{\phi(p,v)=A_p b(v)},\qquad A_p\in\R^{n\times N_c}.
 \label{eq:map}
\end{equation}
For ordinary ODMA, $A_p=S_p/\sqrt{N_c}$, where $S_p$ inserts $N_c$ symbols at distinct channel positions:
$S_p^T S_p=I$. Our earlier notation $\phi=W_p x$, $x=b/\sqrt{N_c}$, is identical with $A_p=W_p/\sqrt{N_c}$.
The matrices $S_p$ specify placement; the polar code specifies which sign vectors are valid. They are separate choices.

\subsection{What the XOR transform actually does}
One kernel sends $(u_0,u_1)$ to $(u_0\oplus u_1,u_1)$, with $\oplus$ denoting addition modulo two.
Applying this recursively gives the row-vector convention $c^T=u^T F^{\otimes q}$,
$F=\left[\begin{smallmatrix}1&0\\1&1\end{smallmatrix}\right]$, $N_c=2^q$.
For example, $u=(0,0,a,b)$ gives $c=(a\oplus b,b,a\oplus b,b)$; $(a,b)=(1,0)$ gives signs $(-1,+1,-1,+1)$.
The transform itself is invertible and adds no redundancy. Fixing input bits and imposing CRC constraints restricts
the valid outputs to $2^{B_c}$ possibilities. Encoding costs $O(N_c\log N_c)$, without enumerating those possibilities.
This is the structural starting point of polar coding~\cite{arikan}; our code uses a mother-code reliability ordering,
not the full 5G NR transport chain.

\clearpage
\section{D2: derive the implemented update}
D2 represents uncertainty about \emph{every complete message}, not individual polar bits.
At iteration $t$, $a_j^t\in[0,1]$ estimates activity and $v_j^t=a_j^t(1-a_j^t)$ is its Bernoulli variance.
Initially $a_j^0=\rho=K/M$. Define
\begin{equation}
 G=\Phi^T\Phi,\qquad d_j=G_{jj},\qquad s_j=\sum_\ell G_{j\ell}^2.
\end{equation}
Although useful in the derivation, the $M\times M$ matrix $G$ is not built in the implementation.

\subsection{Exact interference-cancellation identity}
Form the residual $r^t=y-\Phi a^t$, then take its adjoint correlations $g^t=\Phi^T r^t$.
Adding back the current coordinate's own contribution gives
\begin{align}
 \zeta_j^t&=g_j^t+d_j a_j^t\\
 &=d_j a_j^\star+\sum_{\ell\ne j}G_{j\ell}(a_\ell^\star-a_\ell^t)+\phi_j^T z.
 \label{eq:exactd2}
\end{align}
This identity is exact. The remaining term combines physical noise and errors in cancelling other messages.
It motivates a scalar observation model $\zeta_j^t\approx d_j a_j^\star+\mathcal N(0,\tau_j^t)$.

\subsection{Where approximation enters}
If cancellation errors were uncorrelated, their interference variance would be
$\sum_{\ell\ne j}G_{j\ell}^2 v_\ell^t$. D2 replaces these individual variances by their mean:
\begin{equation}
 \bar v_{-j}^t=\frac{\sum_\ell v_\ell^t-v_j^t}{M-1},\qquad
 \tau_j^t=\softplus(\beta_t)
 \left[\underbrace{\sigma^2d_j}_{\text{physical noise}}+
 \underbrace{\bar v_{-j}^t(s_j-d_j^2)}_{\text{residual interference}}\right].
 \label{eq:tau}
\end{equation}
Equal error variances would justify this averaging \emph{within the uncorrelated-error model}. Iterative reuse of $y$
and the fixed-$K$ constraint generally violate independence. The learned scale calibrates the combined approximation;
it does not separately identify physical noise and interference. There is no Onsager correction or state-evolution
guarantee: this is not an AMP derivation.

\subsection{Gaussian evidence, cardinality and damping}
For Gaussian hypotheses with means $d_j$ and $0$, the log likelihood ratio is
\begin{equation}
 e_j^t=\log\frac{p(\zeta_j^t\mid a_j^\star=1)}{p(\zeta_j^t\mid a_j^\star=0)}
 =\frac{d_j(\zeta_j^t-d_j/2)}{\tau_j^t}.
\end{equation}
D2 then uses
\begin{equation}
 q_j^t=\sigm\!\left(\log\frac{\rho}{1-\rho}+e_j^t+\lambda_t\right),\quad
 \sum_j q_j^t=K,\qquad
 a^{t+1}=\delta_t a^t+(1-\delta_t)q^t,
 \quad\delta_t=\sigm(\gamma_t).
\end{equation}
The scalar $\lambda_t$ enforces total \emph{mean} activity. This is not the exact posterior obtained by conditioning
independent Bernoulli variables on their sum. There are two learned scalars per layer, $(\beta_t,\gamma_t)$: 16 for
the eight-layer experiment. The final readout projects onto nonnegative integer counts summing to $K$.

Training backpropagates through the residual, geometry, scalar inference and encoder. The existing objective combines
layerwise balanced activity BCE on $e^t$ with a small count-regression term; it is not likelihood maximization.
IID evaluation additionally tests a mismatch: repeated messages can have count greater than one, whereas D2's
internal prior remains Bernoulli. None of these approximations makes D2 an optimal decoder.

\clearpage
\section{Which matrix operations can polar structure replace exactly?}
Collect the signs into $\mathcal B=[b(v)]_v\in\R^{N_c\times V}$. The explicit matrix is
\begin{equation}
 \Phi=[A_0\mathcal B\mid A_1\mathcal B\mid\cdots\mid A_{P-1}\mathcal B].
\end{equation}
Partition $a=(a_0,\ldots,a_{P-1})$ with $a_p\in\R^V$. Then
\begin{equation}
 \Phi a=\sum_p A_p m_p,\quad m_p=\mathcal B a_p\in\R^{N_c},\qquad
 (\Phi^T r)_{p,v}=b(v)^T A_p^T r.
 \label{eq:factor}
\end{equation}
These are exact factorizations, not alternative decoders. The first needs only each pattern's mean symbol vector
\emph{after that mean has been computed}; the second turns message correlation into a score on a valid polar word.
But arbitrary $a_p$ still has $V$ entries. A mean symbol vector does not determine the nonlinear next D2 update,
its message-specific variance or its best messages. There is no exact closed update on $m_p$ established here.

\subsection{An exact shortcut for D2's geometry}
Let $\ell_i^T$ be row $i$ of the binary generator $L$. Over all equally weighted payloads,
\begin{equation}
 \frac1V\sum_v b_i(v)b_j(v)
 =\frac1V\sum_v(-1)^{(\ell_i\oplus\ell_j)^T v}
 =\mathbf 1\{\ell_i=\ell_j\}=:Q_{ij}.
 \label{eq:character}
\end{equation}
If the two binary rows differ, pairing each $v$ with a payload that flips their parity cancels the sum.
If they agree, every summand is one. Thus $\mathcal B\mathcal B^T=VQ$ \emph{exactly}.
Repeated rows, including repeated zero rows, matter: assuming $Q=I$ would be wrong in general.

Define the small frame operator $\mathcal F=\Phi\Phi^T\in\R^{n\times n}$. Then
\begin{equation}
 \boxed{\mathcal F=V\sum_p A_p Q A_p^T},\qquad
 d_{p,v}=\|A_p b(v)\|^2,\qquad
 s_{p,v}=b(v)^T A_p^T\mathcal F A_p b(v).
 \label{eq:frame}
\end{equation}
This computes the frame without summing over all payloads or forming $G$. Job 034 uses this identity, with exact
factorized forward/adjoint products. Tests compare both values \emph{and gradients} against explicit multiplication
and the original D2 geometry. The resulting D2 recurrence is unchanged, up to floating-point arithmetic.

\subsection{A useful bridge, not a solved inverse problem}
For a fixed pattern, maximizing the initial correlation is
\begin{equation}
 \argmaxop_{v\text{ valid}} b(v)^T t_p,\qquad t_p=A_p^T r.
\end{equation}
This has the form of decoding a polar word from bit likelihoods. For white Gaussian noise and constant-energy
codewords it is also the isolated-message ML score, with bit LLRs $2t_p/\sigma^2$.
A bounded polar list could therefore propose messages without scanning all $V$.
That is a credible computational bridge, not an exact replacement for D2: later D2 evidence contains
message-dependent uncertainty and self-estimates; list pruning can omit true messages; unknown patterns and
multiple messages on one pattern remain. A compact \emph{atom generator} is not yet a compact \emph{search algorithm}.

\textbf{Scaling boundary.} Per-message polar encoding is polynomial in $N_c$; storing a bank of $P$ patterns is
exponential only in $B_p$, not the full $B$. Our certification adapter still stores local payload signs and maintains
$M=PV$ decoder states. It is intentionally restricted to small $B$. The geometry identity alone does not change this.

\clearpage
\section{What a polar decoder supplies}
Suppose a front end supplies likelihood information about the coded bits. Write its LLR as
\[
\ell_i=\log\frac{p(\text{observation}_i\mid c_i=0)}{p(\text{observation}_i\mid c_i=1)}.
\]
With conditionally independent bit observations and a uniform prior over valid codewords,
\begin{equation}
 P(c\mid\ell)\ \propto\ \mathbf1\{c\text{ valid}\}
 \exp\!\left(\frac12\sum_i\ell_i(-1)^{c_i}\right).
 \label{eq:polarpost}
\end{equation}
The validity indicator contains all the polar/frozen/CRC constraints. The expensive operation is summing over valid
words to obtain marginals. Different polar decoders approximate or organize this inference differently.

\subsection{Successive cancellation and lists}
For the two-bit kernel, independent incoming LLRs combine as
\begin{align}
 f(\ell_1,\ell_2)&=2\operatorname{atanh}\!\left(\tanh(\ell_1/2)\tanh(\ell_2/2)\right),\\
 g(\ell_1,\ell_2,\widehat u_0)&=(-1)^{\widehat u_0}\ell_1+\ell_2.
\end{align}
The $f$ operation infers the XOR bit by marginalizing the other bit. Once that XOR is hypothesized, $g$ aligns
the two pieces of evidence for the second bit. Recursive application gives successive cancellation (SC): frozen
bits are fixed, information bits are decided sequentially. A wrong early decision can propagate.

SCL branches on both information-bit values and retains a bounded number $Q_{\rm list}$ of paths, ranked by their
likelihood metrics~\cite{tal}. CRC filters completed candidates. Standard efficient implementations have complexity
$O(Q_{\rm list}N_c\log N_c)$ rather than $O(2^{B_c})$. A finite list is not exact ML, and passing CRC does not prove
correctness. Our existing native receiver uses a hard CRC-aided list; it is not a soft-output joint receiver.

\subsection{Soft output and the information sent back}
A soft-output decoder estimates coded-bit posterior LLRs
\begin{equation}
 \ell_i^{\rm APP}=\log\frac{\sum_{c\text{ valid}:c_i=0}P(c\mid\ell)}
 {\sum_{c\text{ valid}:c_i=1}P(c\mid\ell)}.
\end{equation}
Polar BP exchanges messages on the XOR graph. Alternatively, soft-output SCL approximates these sums from a list.
Normalizing only over surviving paths can be severely overconfident when the correct path was pruned.
Yuan et al. explicitly estimate the unvisited probability mass~\cite{softscl}; this remains an approximation, not an
exact posterior or a guarantee of list coverage.

In an iterative detector/decoder loop, return \emph{extrinsic} information:
\begin{equation}
 \ell_i^{\rm ext}=\ell_i^{\rm APP}-\ell_i^{\rm in}.
\end{equation}
Here $\ell_i^{\rm in}$ must include exactly the incoming evidence already supplied by the detector, with any other
priors accounted for consistently. Returning the full APP as independent new evidence would repeatedly count the
same observation. Extrinsic subtraction does not make a loopy graph exact, but specifies the intended separation.

\textbf{What polar structure buys.} It replaces an arbitrary prior on $2^{B_c}$ complete payloads by local algebraic
constraints and bounded recursive inference. It does not turn a superposition of several unknown codewords into
independent bit channels. Producing credible $\ell_i$ from that superposition is the multiuser detector's job.

\clearpage
\section{Where compressed sensing and polar inference meet}
An exact posterior can be written using $K$ exchangeable message slots, not device identities:
\begin{equation}
 P(\{p_k,v_k\}_{k=1}^K\mid y)\ \propto\
 \exp\!\left[-\frac{\|y-\sum_k A_{p_k}b(v_k)\|^2}{2\sigma^2}\right]
 \pi(\{p_k,v_k\}).
 \label{eq:joint}
\end{equation}
The prior $\pi$ must match distinct or IID sampling. Equivalently, expose the coded bits and impose the local polar,
frozen and CRC constraints. Both descriptions represent the same channel. A factor graph reorganizes the problem;
its channel-sum factors still couple users, so it is not a tractability proof.

\subsection{Derive a soft interference step, conditional on patterns}
For clarity, first suppose the $K$ patterns are known, and stack their coded signs into $b$:
$y=H b+z$, $H=[A_{p_1}\mid\cdots\mid A_{p_K}]\in\R^{n\times KN_c}$.
Index a sign by $e$. Incoming extrinsic LLRs give $m_e=\tanh(\ell_e^{\rm ext}/2)$ and $v_e=1-m_e^2$.
A diagonal, Gaussian interference approximation for resource $i$ is
\begin{equation}
 r_{i\to e}=y_i-\sum_{f\ne e}H_{if}m_f,\qquad
 T_{i\to e}=\sigma^2+\sum_{f\ne e}H_{if}^2v_f.
\end{equation}
The two hypotheses now have means $+H_{ie}$ and $-H_{ie}$, hence
\begin{equation}
 \boxed{\ell_{i\to e}=\frac{2H_{ie}r_{i\to e}}{T_{i\to e}}}.
 \label{eq:softic}
\end{equation}
Under a factorized resource approximation, sum these messages into each sign's detector LLR and send them to its
polar decoder. Return extrinsic means and uncertainties, then repeat. This is the same basic cancellation/effective
noise/Bayesian-update logic as D2, but the latent variables are coded signs rather than complete-message activities.

A vector version keeps colored interference: $R_{\to e}=\sigma^2I+\sum_{f\ne e}v_fh_fh_f^T$,
$r_{\to e}=y-\sum_{f\ne e}h_fm_f$, giving $\ell_{\to e}=2h_e^TR_{\to e}^{-1}r_{\to e}$.
This motivates soft LMMSE front ends. Both versions still approximate correlated code-bit uncertainty and the effect
of reusing observations. Correlation-aware updates cost more; adding a learned scale does not make either exact.

\begin{center}
\begin{tabularx}{\textwidth}{@{}lYY@{}}\toprule
 & Full-message D2 & Structured detector + polar inference\\\midrule
Latent state & $M$ message activities & Pattern/activity state and coded-bit messages\\
Linear evidence & $\Phi^T(y-\Phi a)$ & Placement/spreading residuals and bit LLRs\\
Prior step & Bernoulli + mean cardinality & Polar/CRC constraints + activity model\\
Main limitation & Exponential state; approximate uncertainty & Unknown-pattern search, collisions, loopy uncertainty\\\bottomrule
\end{tabularx}
\end{center}

\textbf{The unresolved part is substantial.} Known patterns above are an explanatory conditioning, not an allowed
oracle in the final receiver. Different full messages can share a pattern. With IID prefixes, a tagged user's pattern
collides with probability $1-(1-1/P)^{K-1}$, even when full-message collisions are negligible. A one-packet-per-pattern
model would therefore change the problem. Activity/pattern inference must allow multiplicity, or a revisable list must
represent competing messages and measure its misses. Simply replacing D2's sigmoid by SC does neither.

\clearpage
\section{Existing work already covers important parts of this overlap}
This is a focused map of relevant primary papers, not an exhaustive survey or a claim that their assumptions match
Job 033. The opportunity must be narrower than ``combine polar coding with interference cancellation''.

\textbf{Polar coding plus spreading in URA.} Pradhan et al.~\cite{spreading} split the message into a spreading
signature and a polar-coded payload. Their receiver detects signatures, forms LLRs with MMSE estimation, applies
CRC-aided SCL, and cancels recovered packets. Thus polar coding without ODMA is already an established construction.
It is a natural architectural reference, not a new idea in itself. Their large-payload operating point is not our
short-message certification regime.

\textbf{A soft-in/soft-out multiuser loop in URA.} Pradhan et al.~\cite{ldpc} replace the polar/list stage of that
spreading architecture by LDPC BP and exchange soft information with an MMSE detector. This directly addresses
the value of revisable uncertainty instead of only hard cancellation. It also shows that the code and the iterative
detector can be designed together; polar is not the only possible constraint engine.

\textbf{Polar BP with soft interference cancellation.} Ebada et al.~\cite{polarsoic} study iterative polar decoding
in two- and four-user Gaussian MACs using an IDMA construction. This is a close reference for the bit-level loop in
Section 5, but it does not solve unsourced pattern/activity search. Its complexity claims should not be imported
as a complete URA decoder guarantee.

\textbf{Soft output from a strong list decoder.} Yuan, Duffy and M\'edard~\cite{softscl} supply a route to calibrated
soft information from SCL rather than requiring polar BP. This is relevant if plain BP gives up too much decoding
quality. Their treatment of probability outside the list is particularly pertinent to candidate-based inference.

\textbf{Learning the polar decoder, not only the code.} Xu et al.~\cite{learnedpolar} learn weighted message updates
on a sparse polar decoding graph. Learned decoding of fixed polar constraints is therefore established prior art.
A project-specific question would be whether learned multiuser uncertainty exchange helps, under matched waveforms,
activity assumptions and computational budgets; it is not merely whether BP weights can be trained.

\textbf{ODMA and soft polar feedback already meet.} Li et al.~\cite{mimoodma} develop a MIMO URA architecture with
pilots, hierarchical pattern detection and soft polar feedback. It is directly relevant to joint pattern/data
inference. Its multi-antenna channel estimation and pilot structure differ from our real, known-channel MAC, so it
is neither an interchangeable baseline nor evidence that our remaining inverse problem has already been solved.

\subsection{What seems worth deriving next}
The most informative comparison is between two uncertainty representations on \emph{the same} fixed encoder:
\begin{enumerate}
\item A revisable polar-generated codeword list with joint residual refinement. Separate proposal recall and
      list size from the error of refinement; do not supply true messages or hide misses.
\item A soft detector/code loop, with explicit activity and same-pattern collision handling. Compare scalar
      interference variances with a colored/LMMSE alternative before adding many learned parameters.
\end{enumerate}
Use known-pattern tests only to diagnose the inference core, then remove that oracle. Keep known $K$ initially for
comparison with D2, but label it; estimating $K$ and controlling false alarms remain separate requirements.

\textbf{Research position.} D2's gain suggests useful uncertainty is lost by the current native receiver, not that
D2's equations are uniquely correct. A tractable decoder may improve on D2 by using code constraints more faithfully.
Conversely, better single-user polar decoding alone may not repair poor pattern evidence or an overconfident front end.
No receiver choice is made by the encoder tests on the next page.

\clearpage
\section{Job 034: two controlled encoder investigations}
Keep the selected Job 033 polar/CRC code, pattern placements, seeds, observations and D2 training contract.
Vary one part of $A_p$ at a time. The unchanged ODMA--polar/D2 rows from 033 supply the fixed baseline;
dense/sparse D2 and native ODMA rows from 032 are reused rather than retrained.

\subsection{Investigation 2: learn symbol-energy profiles}
For each pattern, learn $\theta_p\in\R^{N_c}$ and set
\begin{equation}
 e_{p,i}=\frac{\exp\theta_{p,i}}{\sum_j\exp\theta_{p,j}},\quad
 A_p=S_p\diag(\sqrt{e_p}),\qquad
 \|\phi(p,v)\|^2=\sum_i e_{p,i}=1\quad\text{for every }v.
\end{equation}
Initialize $\theta=0$ to recover equal energy exactly. Polar bits, CRC and positions are fixed; the continuous energy
logits and D2 parameters are trained jointly. This changes the relative importance of coded symbols and may change
within-pattern distances. It is a structured $P N_c$-parameter family, not an arbitrary amplitude per message.
The normalization is enforced in every forward call, including for previously unseen payloads.

\subsection{Investigation 3: learn one shared pre-placement mixer}
Learn a skew-symmetric matrix $K_0$ through its $N_c(N_c-1)/2$ upper-triangular entries and use the Cayley map
\begin{equation}
 U=(I-K_0)^{-1}(I+K_0),\quad U^TU=I,\qquad
 A_p=\frac1{\sqrt{N_c}}S_p U.
\end{equation}
Initialize $K_0=0$, giving $U=I$ and exactly the original codebook. This finite chart covers orthogonal matrices
without eigenvalue $-1$ in the identity component; it is not an unrestricted real mixer.
It preserves unit energy and every \emph{within-pattern} distance, while changing cross-pattern interactions:
\begin{equation}
 \phi(p,v)^T\phi(q,v')=\frac1{N_c}b(v)^T U^T S_p^T S_q U b(v').
\end{equation}
It is not a common rotation of the entire received waveform. An isolated known-pattern signal can be unmixed by
$U^T S_p^T$, but multiuser interference becomes mixed too. Dense $U$ adds $O(N_c^2)$ per-message work and can complicate
future bit-level inference. Better exhaustive D2 performance would not automatically imply a practical receiver gain.

\subsection{Contract, checks and interpretation}
\begin{itemize}
\item Twelve rows: two families $\times$ $B=12,14$ $\times$ three paired seeds; $n=256$.
      Selected $N_c=32$, $B_p=6,8$: energy has 2,048/8,192 parameters; the shared mixer has 496.
\item Eight-layer D2, 200 epochs maximum, patience 10 and restoration of the best encoder/decoder pair.
      Same loads $7,15,22,26$, physical $E_b/N_0=-4{:}2{:}10$ dB evaluation, distinct-message training and distinct/IID
      evaluation as 033. No energy-plus-mixer combined arm and no new receiver.
\item Exact initialization matches all six returned 033 codebooks. Forward/adjoint products, geometry and gradients
      match explicit references; all generated codewords retain unit energy. A sampled $B=100$ forward check tests
      generation only. The small-$B$ D2 adapter explicitly rejects that decoding claim.
\end{itemize}
Local $B=8$ mini-runs reduced validation loss from about 0.1082 to 0.0821 (fixed), 0.0758 (energy), and 0.0741 (mixer).
These are execution and initial-learning checks, not production rankings; lower surrogate loss need not give lower
PUPE in a tiny test. HPC results are pending. Compare paired per-load/SNR gaps against fixed polar D2 and dense D2:
a gain identifies useful waveform freedom; no gain at this budget is not a proof of optimality.

\clearpage
\section*{References and implementation pointers}
The derivation follows the real, unit-energy comparison system. It does not import fading, pilot or antenna gains
from papers using different models. Report equations were checked against the retained implementations.

\begin{thebibliography}{9}
\setlength{\itemsep}{0.5em}
\bibitem{arikan} E. Ar\i kan, \emph{Channel Polarization: A Method for Constructing Capacity-Achieving Codes for
Symmetric Binary-Input Memoryless Channels}, 2009. \href{https://arxiv.org/abs/0807.3917}{arXiv:0807.3917}.
\bibitem{tal} I. Tal and A. Vardy, \emph{List Decoding of Polar Codes}, 2015 (preprint 2012).
\href{https://arxiv.org/abs/1206.0050}{arXiv:1206.0050}.
\bibitem{spreading} A. K. Pradhan, V. K. Amalladinne, K. R. Narayanan and J.-F. Chamberland,
\emph{Polar Coding and Random Spreading for Unsourced Multiple Access}, 2019 preprint / ICC 2020.
\href{https://arxiv.org/abs/1911.01009}{arXiv:1911.01009}.
\bibitem{ldpc} A. K. Pradhan, V. K. Amalladinne, K. R. Narayanan and J.-F. Chamberland,
\emph{LDPC Codes with Soft Interference Cancellation for Uncoordinated Unsourced Multiple Access}, 2021.
\href{https://arxiv.org/abs/2105.13985}{arXiv:2105.13985}.
\bibitem{polarsoic} M. Ebada, S. Cammerer, A. Elkelesh, M. Geiselhart and S. ten Brink,
\emph{Iterative Detection and Decoding of Finite-Length Polar Codes in Gaussian Multiple Access Channels}, 2020.
\href{https://arxiv.org/abs/2012.01075}{arXiv:2012.01075}.
\bibitem{softscl} P. Yuan, K. R. Duffy and M. M\'edard, \emph{Soft-Output Successive Cancellation List Decoding},
IEEE Transactions on Information Theory, 71(2), 1007--1017, 2025.
\href{https://arxiv.org/abs/2412.03623}{arXiv:2412.03623}.
\bibitem{learnedpolar} W. Xu, X. You, C. Zhang and Y. Be'ery,
\emph{Polar Decoding on Sparse Graphs with Deep Learning}, 2018.
\href{https://arxiv.org/abs/1811.09801}{arXiv:1811.09801}.
\bibitem{mimoodma} T. Li et al., \emph{ODMA-based MIMO Massive Unsourced Random Access with Soft-Output Polar Codes},
2026 preprint; listed as accepted in IEEE Transactions on Wireless Communications.
\href{https://arxiv.org/abs/2608.28085}{arXiv:2608.28085}.
\end{thebibliography}

\textbf{Implemented sources.}
\begin{itemize}
\item \texttt{framework/learned\_decoders.py}: unchanged D2 recurrence and its approximations.
\item \texttt{baselines/polar.py}, \texttt{baselines/odma\_polar.py}: binary code, placement and native list/SIC receiver.
\item \texttt{framework/polar\_mapping.py}: procedural mappings, small-$B$ adapter and exact geometry shortcut.
\item \texttt{benchmarks/polar\_mapping.py}: unchanged training/evaluation contract via the shared benchmark.
\item \texttt{jobs/034\_polar\_mapping/README.md}: manifest, controls, commands and local checks.
\item \texttt{tests/polar\_mapping\_test.py}: algebra/gradient regressions;
      \texttt{tests/polar\_mapping\_learning.py}: smoke and mini runs.
\end{itemize}
\end{document}
